% !TEX root = ./BaoCao_AI_Agent_Platform.tex
% !TEX program = pdflatex
% !TEX encoding = UTF-8

\documentclass[12pt,a4paper]{report} 
\usepackage[utf8]{inputenc}
\usepackage[vietnamese]{babel}
\usepackage{amsmath, amsfonts, amssymb}
\usepackage{graphicx}
\usepackage{hyperref}
\usepackage{xurl}
\usepackage[protrusion=false]{microtype}
\usepackage{geometry}
\usepackage{listings} 
\usepackage{color}
\usepackage{indentfirst} 
\usepackage{tikz}
\usetikzlibrary{shapes.geometric, arrows.meta, positioning, fit, calc}
\usepackage{float}
\usepackage{enumitem}

% --- CẤU HÌNH SIÊU LIÊN KẾT VÀ TỰ ĐỘNG NGẮT DÒNG URL ---
\hypersetup{
    colorlinks=true,
    linkcolor=black,
    citecolor=black,
    urlcolor=blue,
    breaklinks=true
}

% --- CẤU HÌNH CHỐNG TRÀN LỀ PHẢI (TRÁNH OVERFULL HBOX) ---
\emergencystretch 3em
\tolerance=2000
\hbadness=2000
\sloppy

% --- CẤU HÌNH FONT CHỮ TIMES NEW ROMAN ---
\usepackage{mathptmx}

% --- CẤU HÌNH CĂN LỀ THEO CHUẨN ---
% Trái 2.5cm; Trên, dưới 2.5cm; Phải 2.0cm
\geometry{top=2.5cm, bottom=2.5cm, left=2.5cm, right=2cm}

% --- CẤU HÌNH DÃN DÒNG VÀ KHOẢNG CÁCH ĐOẠN ---
\usepackage{setspace}
\setstretch{1.3}
\setlength{\parskip}{6pt}
\setlength{\parindent}{1.27cm}

% --- CẤU HÌNH DANH SÁCH (GẠCH ĐẦU DÒNG VÀ ĐÁNH SỐ CHUẨN GỌN) ---
\setlist[itemize]{
    label={-},
    leftmargin=1.27cm,
    labelsep=0.5em,
    itemsep=2pt,
    parsep=0pt,
    topsep=3pt,
    partopsep=0pt
}
\setlist[enumerate]{
    leftmargin=1.27cm,
    labelsep=0.5em,
    itemsep=2pt,
    parsep=0pt,
    topsep=3pt,
    partopsep=0pt
}

% --- CẤU HÌNH HEADER VÀ FOOTER ---
\usepackage{fancyhdr}
\pagestyle{fancy}
\fancyhf{} 
\lhead{Báo cáo: Nghiên cứu và xây dựng AI Agent hỗ trợ Code Review sử dụng Spring AI} 
\rfoot{\thepage} 
\renewcommand{\headrulewidth}{0.4pt}
\renewcommand{\footrulewidth}{0pt}

% --- CẤU HÌNH HEADING (STYLE CHUẨN) ---
\usepackage{titlesec}

% Heading 1 (Chương 1, Chương 2...): Size 18pt, Bold
\titleformat{\chapter}[hang]
  {\fontsize{18}{22}\bfseries\mathversion{bold}}
  {Chương \thechapter.}
  {0.5em}
  {}
\titlespacing*{\chapter}{0pt}{-20pt}{20pt} 

% Heading 2 (1.1, 1.2...): Size 13pt, Bold
\titleformat{\section}
  {\fontsize{13}{16}\bfseries\mathversion{bold}}
  {\thesection.}
  {0.5em}
  {}

% Heading 3 (1.1.1, 1.1.2...): Size 13pt, Italic
\titleformat{\subsection}
  {\fontsize{13}{16}\itshape}
  {\thesubsection.}
  {0.5em}
  {}

\titleformat{\subsubsection}
  {\fontsize{13}{16}\itshape}
  {\thesubsubsection.}
  {0.5em}
  {}

% --- CẤU HÌNH CAPTION ---
\usepackage{caption}
\captionsetup{font={small,it}, labelsep=colon}

% --- CẤU HÌNH HÌNH HIỂN THỊ CODE ---
\definecolor{codegreen}{rgb}{0,0.6,0}
\definecolor{codegray}{rgb}{0.5,0.5,0.5}
\lstset{
    commentstyle=\color{codegreen},
    keywordstyle=\color{blue},
    numberstyle=\tiny\color{codegray},
    stringstyle=\color{red},
    basicstyle=\ttfamily\small,
    breaklines=true,
    captionpos=b,
    keepspaces=true,
    numbers=left,
    showspaces=false,
    inputencoding=utf8,
    extendedchars=true
}

\setlength{\headheight}{15pt}
\addtolength{\topmargin}{-3pt}

\begin{document}

% ==============================================
% TRANG BÌA
% ==============================================
\begin{titlepage}
    \centering 
    
    {\fontsize{14pt}{17pt}\selectfont \bfseries TRƯỜNG ĐẠI HỌC KHOA HỌC TỰ NHIÊN\par}
    \vspace{0.25cm}
    {\fontsize{14pt}{17pt}\selectfont \bfseries ĐẠI HỌC QUỐC GIA HÀ NỘI\par}
    
    \vfill 
    \vspace{0.5cm}
    
    \vspace{0.75cm}
    
    {\fontsize{16pt}{19pt}\selectfont \bfseries BÁO CÁO THỰC TẬP HÈ\par}
    \vspace{0.25cm}
    
    {\fontsize{20pt}{24pt}\selectfont \bfseries NGÀNH KHOA HỌC MÁY TÍNH \& THÔNG TIN\par} 
    
    \vspace{0.75cm}
    
    {\fontsize{14pt}{17pt}\selectfont \bfseries \raggedright Chủ đề:\par}
    \vspace{0.25cm}
    
    {\fontsize{18pt}{22pt}\selectfont \bfseries NGHIÊN CỨU VÀ XÂY DỰNG AI AGENT HỖ TRỢ CODE REVIEW SỬ DỤNG SPRING AI\par}
    
    \vfill
    
    \vspace{0.5cm}
    \begin{minipage}{0.9\textwidth} 
        \raggedright 

        {\fontsize{12pt}{15pt}\selectfont \bfseries Cán bộ hướng dẫn DN: Kỹ sư Trần Văn C \par}
        \vspace{0.25cm}
        
        {\fontsize{12pt}{15pt}\selectfont \bfseries Sinh viên thực hiện: Nguyễn Việt Anh\par}
        \vspace{0.25cm}
        
    \end{minipage}
    
    \vspace{1cm} 
    
    {\fontsize{12pt}{15pt}\selectfont \bfseries Hà Nội, Năm 2026\par}

    \begin{tikzpicture}[remember picture, overlay]
        \draw[line width=3pt, rounded corners=0pt, color=black] 
            ([shift={(1cm,-1cm)}]current page.north west) rectangle
            ([shift={(-1cm,1cm)}]current page.south east);
    \end{tikzpicture}
    
\end{titlepage}

% ==============================================
% LỜI MỞ ĐẦU & MỤC LỤC
% ==============================================
\phantomsection

\chapter*{Lời mở đầu}
\addcontentsline{toc}{chapter}{Lời mở đầu}

Trong quy trình phát triển phần mềm hiện đại, việc đảm bảo chất lượng, tính ổn định và mức độ an toàn bảo mật của mã nguồn luôn là một trong những ưu tiên hàng đầu của các tổ chức công nghệ. Hoạt động đánh giá mã nguồn (Code Review) đóng vai trò như một mắt xích kiểm thử then chốt trước khi mã nguồn được tích hợp vào nhánh chính và đưa lên môi trường vận hành. Tuy nhiên, phương pháp đánh giá thủ công truyền thống do các kỹ sư phần mềm đảm nhiệm thường đòi hỏi nhiều thời gian, tạo điểm nghẽn trong tiến độ phát hành tính năng và phụ thuộc lớn vào sự tập trung cũng như kinh nghiệm cá nhân của người đánh giá.

Sự phát triển mạnh mẽ của Trí tuệ nhân tạo, đặc biệt là các mô hình ngôn ngữ lớn kết hợp cùng kiến trúc tác tử (AI Agent) và cơ chế gọi công cụ (Tool Calling), đã mở ra hướng tiếp cận mới trong việc tự động hóa quá trình phân tích mã nguồn. Khác với các công cụ phân tích tĩnh truyền thống chỉ kiểm tra theo tập quy tắc cố định, tác tử AI có khả năng chủ động điều hướng, khám phá cấu trúc thư mục dự án và đưa ra lời giải thích cùng các khuyến nghị khắc phục bằng ngôn ngữ tự nhiên.

Xuất phát từ yêu cầu thực tiễn trong đợt thực tập tốt nghiệp, em đã thực hiện đề tài: \textbf{``Nghiên cứu và xây dựng AI Agent hỗ trợ Code Review sử dụng Spring AI''}. Báo cáo này tổng hợp quá trình tìm hiểu lý thuyết, phân tích thiết kế và hiện thực hóa một hệ thống dịch vụ backend trên nền tảng Spring Boot và Java 21, tích hợp thư viện Spring AI để điều khiển mô hình Google Gemini thực thi phân tích mã nguồn, kết hợp cơ chế xử lý bất đồng bộ, các giải pháp bảo mật đa lớp và kiểm thử tự động đo lường độ phủ với JaCoCo.

Bố cục của báo cáo được tổ chức thành 5 chương chính:
\begin{itemize}
    \item \textbf{Chương 1: Giới thiệu bài toán và cơ sở công nghệ} -- Trình bày bối cảnh, thách thức của bài toán đánh giá mã nguồn, so sánh các phương pháp tiếp cận và cơ sở lý thuyết về AI Agent, Spring AI.
    \item \textbf{Chương 2: Phân tích và thiết kế hệ thống} -- Đặc tả các yêu cầu chức năng, phi chức năng, ma trận phân quyền RBAC, sơ đồ kiến trúc phân tầng, cơ sở dữ liệu, luồng xử lý bất đồng bộ và thiết kế RESTful API.
    \item \textbf{Chương 3: Quá trình triển khai và hiện thực hóa hệ thống} -- Trình bày chi tiết việc cài đặt tác tử AI, kỹ thuật thiết kế lời nhắc hệ thống, quản lý luồng bất đồng bộ \texttt{ThreadPoolTaskExecutor}, bảo mật phòng vệ hệ thống tệp và đóng gói Docker.
    \item \textbf{Chương 4: Kết quả thử nghiệm và đánh giá hệ thống} -- Tổng kết kết quả kiểm thử chức năng qua Postman, số liệu đo độ phủ kiểm thử tự động qua JaCoCo, thực nghiệm trên các dự án mẫu và phân tích sự cố kỹ thuật.
    \item \textbf{Chương 5: Kết luận và hướng phát triển} -- Đánh giá tổng thể các kết quả đạt được, giá trị thực tiễn, bài học kinh nghiệm và các định hướng mở rộng trong tương lai.
\end{itemize}

Em xin bày tỏ lòng biết ơn sâu sắc tới các thầy cô giáo tại Trường Đại học Khoa học Tự nhiên, Đại học Quốc gia Hà Nội đã tận tình truyền đạt những kiến thức chuyên môn quý báu trong suốt quá trình học tập. Em cũng xin chân thành cảm ơn Ban lãnh đạo đơn vị thực tập cùng Cán bộ hướng dẫn tại doanh nghiệp đã tận tâm chỉ bảo, định hướng và tạo mọi điều kiện thuận lợi để em hoàn thành đợt thực tập và xây dựng báo cáo này.

Mặc dù đã có nhiều cố gắng trong quá trình nghiên cứu và triển khai, song do thời gian và kinh nghiệm còn hạn chế, báo cáo khó tránh khỏi những thiếu sót. Em rất mong nhận được những ý kiến đóng góp, chỉ dẫn quý báu từ các thầy cô để đề tài được hoàn thiện hơn nữa.

\vspace{1cm}
\begin{flushright}
Sinh viên thực hiện \\
\textbf{Nguyễn Việt Anh}
\end{flushright}
\newpage

\tableofcontents
\cleardoublepage

\listoffigures
\cleardoublepage

\listoftables
\cleardoublepage

\chapter*{Danh mục từ viết tắt}
\addcontentsline{toc}{chapter}{Danh mục từ viết tắt}

\begin{table}[H]
\centering
\renewcommand{\arraystretch}{1.2}
\begin{tabular}{|p{2.5cm}|p{11.5cm}|}
\hline
\textbf{Từ viết tắt} & \textbf{Diễn giải} \\ \hline
API & Application Programming Interface -- Giao diện lập trình ứng dụng \\ \hline
AST & Abstract Syntax Tree -- Cây cú pháp trừu tượng \\ \hline
CI/CD & Continuous Integration / Continuous Deployment -- Tích hợp và triển khai liên tục \\ \hline
DTO & Data Transfer Object -- Đối tượng truyền tải dữ liệu \\ \hline
ERD & Entity Relationship Diagram -- Sơ đồ thực thể liên kết \\ \hline
IDOR & Insecure Direct Object Reference -- Lỗ hổng tham chiếu đối tượng trực tiếp không an toàn \\ \hline
IoC & Inversion of Control -- Đảo ngược điều khiển \\ \hline
JPA & Java Persistence API -- Giao diện ánh xạ đối tượng quan hệ trong Java \\ \hline
JWT & JSON Web Token -- Mã thông báo định dạng JSON \\ \hline
LLM & Large Language Model -- Mô hình ngôn ngữ lớn \\ \hline
RAG & Retrieval-Augmented Generation -- Kỹ thuật sinh tăng cường truy xuất dữ liệu \\ \hline
RBAC & Role-Based Access Control -- Kiểm soát truy cập dựa trên vai trò \\ \hline
REST & Representational State Transfer -- Phong cách kiến trúc dịch vụ web \\ \hline
SSE & Server-Sent Events -- Cơ chế truyền dữ liệu một chiều từ máy chủ đến máy khách \\ \hline
\end{tabular}
\end{table}
\cleardoublepage

% ==============================================
% NỘI DUNG CHÍNH
% ==============================================
\chapter{GIỚI THIỆU BÀI TOÁN VÀ CƠ SỞ CÔNG NGHỆ}

\section{Bối cảnh và phát biểu bài toán}
\subsection{Bối cảnh thực tế trong quy trình phát triển phần mềm}
Trong các quy trình phát triển phần mềm hiện đại, đánh giá mã nguồn là bước kiểm thử tĩnh quan trọng nhằm rà soát chất lượng, tính ổn định và mức độ an toàn của sản phẩm trước khi mã nguồn được tích hợp vào nhánh chính và đưa lên môi trường vận hành.

Quá trình này giúp:
\begin{itemize}
    \item Phát hiện sớm các lỗi sai trong logic nghiệp vụ, vi phạm cấu trúc thiết kế hoặc các đoạn mã dư thừa.
    \item Hạn chế các nguy cơ bảo mật như để lộ khóa bí mật trong mã nguồn, thiếu kiểm tra quyền truy cập tài nguyên.
    \item Duy trì tính nhất quán về phong cách lập trình, tạo thuận lợi cho việc đọc hiểu và bảo trì dự án.
\end{itemize}

\subsection{Phát biểu bài toán và những thách thức cốt lõi}
Trong thực tế phát triển phần mềm, phương pháp đánh giá mã nguồn thủ công do lập trình viên đảm nhiệm thường gặp một số khó khăn:
\begin{enumerate}
    \item \textbf{Độ trễ thời gian}: Việc đọc hiểu từng dòng thay đổi trên các yêu cầu gộp mã (Pull Request) đòi hỏi nhiều thời gian, dễ gây chậm tiến độ tích hợp tính năng mới.
    \item \textbf{Tính phụ thuộc vào người đánh giá}: Kết quả rà soát phụ thuộc vào kinh nghiệm, sự tập trung và kiến thức bảo mật của từng cá nhân. Một số vấn đề logic hoặc hiệu năng truy vấn cơ sở dữ liệu có thể bị bỏ sót.
    \item \textbf{Khả năng mở rộng}: Khi số lượng lập trình viên và khối lượng mã nguồn tăng lên, việc bố trí nhân sự rà soát thủ công toàn bộ mã nguồn làm gia tăng chi phí vận hành của nhóm phát triển.
\end{enumerate}

Từ thực tế trên, việc nghiên cứu xây dựng giải pháp hỗ trợ đánh giá mã nguồn tự động dựa trên trí tuệ nhân tạo có khả năng thu thập ngữ cảnh dự án và hỗ trợ lập trình viên sơ lọc các vấn đề kỹ thuật là một hướng tiếp cận có giá trị thực tiễn.

\section{Mục tiêu và phạm vi đề tài trong đợt thực tập}
Đề tài tập trung vào các mục tiêu sau:
\begin{itemize}
    \item \textbf{Về nghiên cứu lý thuyết}: Tìm hiểu nguyên lý hoạt động của kiến trúc tác tử AI, mô hình suy luận ReAct và cơ chế gọi công cụ của các mô hình ngôn ngữ lớn.
    \item \textbf{Về xây dựng hệ thống}:
    \begin{itemize}
        \item Xây dựng hệ thống backend trên nền tảng Spring Boot và Java 21 theo kiến trúc phân tầng (Controller, Service, Repository, DTO).
        \item Sử dụng thư viện Spring AI kết nối với mô hình Google Gemini để thực hiện phân tích mã nguồn theo các định hướng: bảo mật, hiệu năng, lỗi logic và chuẩn mã nguồn.
        \item Cài đặt cơ chế xử lý bất đồng bộ với \texttt{ThreadPoolTaskExecutor}, cho phép tác vụ phân tích chạy ở luồng nền để không làm gián đoạn phản hồi của API.
        \item Thiết lập các cơ chế kiểm soát an toàn: xác thực JWT, kiểm tra quyền sở hữu tài nguyên nhằm hạn chế IDOR và chuẩn hóa đường dẫn nhằm ngăn ngừa tấn công duyệt thư mục trái phép (Path Traversal).
    \end{itemize}
    \item \textbf{Về kiểm thử và đóng gói}: Xây dựng bộ kiểm thử tự động đo độ phủ mã nguồn với JaCoCo và cấu hình đóng gói triển khai ứng dụng cùng cơ sở dữ liệu MySQL bằng Docker Compose.
\end{itemize}

\section{Tổng quan và so sánh các phương pháp đánh giá mã nguồn}
Bảng \ref{tab:method_comparison} tổng hợp so sánh giữa 4 phương pháp đánh giá mã nguồn phổ biến hiện nay.

\begin{table}[H]
\centering
\caption{So sánh các phương pháp đánh giá mã nguồn}
\label{tab:method_comparison}
\renewcommand{\arraystretch}{1.2}
\begin{tabular}{|p{3.2cm}|p{2.6cm}|p{2.8cm}|p{2.6cm}|p{2.8cm}|}
\hline
\textbf{Tiêu chí} & \textbf{Thủ công} & \textbf{Phân tích tĩnh} & \textbf{Mô hình RAG} & \textbf{AI Agent Tool Calling} \\ \hline
Hiểu ngữ cảnh dự án & Tốt (dựa vào con người) & Thấp (chỉ quét theo quy tắc cục bộ) & Trung bình (dựa trên tương đồng vector) & Tốt (chủ động khám phá cấu trúc tệp) \\ \hline
Khả năng giải thích & Chi tiết bằng ngôn ngữ tự nhiên & Khô cứng (xuất mã lỗi định sẵn) & Tốt (sinh văn bản tự nhiên) & Tốt (nêu nguyên nhân và gợi ý khắc phục) \\ \hline
Chi phí tài nguyên & Cao (chi phí nhân lực) & Thấp (chạy trực tiếp trên máy chủ) & Đáng kể (tính toán vector nhúng) & Phù hợp (chỉ đọc các tệp trọng tâm) \\ \hline
Độ linh hoạt & Cao & Thấp (phụ thuộc tập quy tắc cố định) & Trung bình & Linh hoạt theo yêu cầu phân tích \\ \hline
Mức độ tự động hóa & Thủ công & Tự động hoàn toàn (CI/CD) & Tự động & Tự động (qua REST API) \\ \hline
\end{tabular}
\end{table}

\subsection{Đánh giá mã nguồn thủ công}
Lập trình viên trực tiếp đọc và nhận xét mã nguồn của nhau trên nền tảng quản lý phiên bản. Phương pháp này có ưu điểm nắm rõ logic nghiệp vụ nhưng tốn thời gian và phụ thuộc vào mức độ tập trung của người đánh giá.

\subsection{Phân tích tĩnh (Static Analysis)}
Các công cụ phân tích tĩnh quét mã nguồn dựa trên cây cú pháp trừu tượng và bộ quy tắc được định nghĩa trước. Phương pháp này có tốc độ thực thi nhanh nhưng hạn chế trong việc hiểu ngữ cảnh tổng thể và khó đưa ra giải thích linh hoạt bằng ngôn ngữ tự nhiên.

\subsection{Mô hình ngôn ngữ lớn kết hợp tìm kiếm tăng cường (RAG)}
Mã nguồn được chia thành các đoạn nhỏ, chuyển đổi thành vector nhúng và lưu vào cơ sở dữ liệu vector. Hạn chế của phương pháp này đối với dự án phần mềm là cấu trúc phân cấp thư mục có thể bị phân mảnh, gây khó khăn cho việc nắm bắt mối quan hệ giữa các tầng kiến trúc.

\subsection{Mô hình AI Agent kết hợp gọi công cụ (Tool Calling)}
Đây là hướng tiếp cận được lựa chọn trong đề tài. AI Agent vận hành theo chu trình suy luận nhiều bước: mô hình chủ động khảo sát cấu trúc thư mục, định vị tệp cấu hình và gọi công cụ đọc nội dung các tệp liên quan khi cần thiết. Cách tiếp cận này giúp duy trì cấu trúc dự án mà không cần lập chỉ mục vector tĩnh từ trước.

\section{Cơ sở lý thuyết và công nghệ sử dụng}

\subsection{Mô hình suy luận ReAct và cơ chế Tool Calling}
Kiến trúc tác tử AI trong đề tài vận hành dựa trên khung suy luận ReAct (Reasoning and Acting):
\begin{enumerate}
    \item \textbf{Quan sát}: Tiếp nhận yêu cầu đánh giá và khảo sát cấu trúc thư mục của dự án.
    \item \textbf{Suy luận}: Mô hình phân tích các tệp tin cần kiểm tra tiếp theo (ví dụ: \texttt{pom.xml} để xem thư viện phụ thuộc, \texttt{application.yml} để xem cấu hình hệ thống).
    \item \textbf{Hành động}: Mô hình sinh yêu cầu gọi công cụ với tên công cụ và tham số tương ứng.
    \item \textbf{Thực thi và phản hồi}: Hệ thống backend thực thi công cụ trên hệ thống tệp cục bộ và gửi trả kết quả để mô hình tiếp tục suy luận cho đến khi hoàn thành báo cáo.
\end{enumerate}

\subsection{Khung làm việc Spring Boot và Spring AI}
\begin{itemize}
    \item \textbf{Spring Boot 3.x / Java 21}: Cung cấp nền tảng phát triển ứng dụng với cơ chế tiêm phụ thuộc (Dependency Injection), quản lý luồng bất đồng bộ và hệ sinh thái bảo mật hoàn chỉnh.
    \item \textbf{Spring AI}: Thư viện trong hệ sinh thái Spring hỗ trợ tích hợp trí tuệ nhân tạo, cung cấp giao diện trừu tượng hóa \texttt{ChatClient} để tương tác với các mô hình ngôn ngữ lớn và hỗ trợ cơ chế đăng ký công cụ tự động.
\end{itemize}

\subsection{Mô hình ngôn ngữ lớn Google Gemini}
Google Gemini được lựa chọn nhờ khả năng tiếp nhận ngữ cảnh lớn, tốc độ phản hồi nhanh và hỗ trợ cơ chế gọi công cụ (Function Calling / Tool Calling).

\subsection{Bảo mật hệ thống với Spring Security và JWT}
Hệ thống sử dụng cơ chế xác thực không trạng thái dựa trên mã thông báo JWT có chữ ký HMAC-SHA256, giúp bảo vệ các điểm cuối API và kiểm soát quyền truy cập tài nguyên của người dùng.

% ==============================================
% CHƯƠNG 2: PHÂN TÍCH VÀ THIẾT KẾ HỆ THỐNG
% ==============================================
\chapter{PHÂN TÍCH VÀ THIẾT KẾ HỆ THỐNG}

\section{Phân tích yêu cầu hệ thống}

\subsection{Yêu cầu chức năng}
Hệ thống được thiết kế đáp ứng 4 nhóm yêu cầu chức năng chính:
\begin{enumerate}
    \item \textbf{Phân hệ xác thực và quản lý tài khoản}:
    \begin{itemize}
        \item Đăng ký tài khoản người dùng với các ràng buộc duy nhất về \texttt{username} và \texttt{email}.
        \item Đăng nhập hệ thống, kiểm tra mật khẩu qua thuật toán băm BCrypt và cấp phát mã thông báo JWT.
    \end{itemize}
    \item \textbf{Phân hệ quản lý dự án mã nguồn}:
    \begin{itemize}
        \item Khai báo thông tin dự án kèm đường dẫn thư mục mã nguồn trên máy chủ.
        \item Tra cứu, cập nhật thông tin và xóa dự án thuộc quyền sở hữu của người dùng.
    \end{itemize}
    \item \textbf{Phân hệ điều phối và đánh giá mã nguồn}:
    \begin{itemize}
        \item Khởi tạo tác vụ phân tích theo các chế độ: \texttt{SECURITY} (bảo mật), \texttt{PERFORMANCE} (hiệu năng), \texttt{BUG} (lỗi logic), \texttt{CODE\_QUALITY} (chuẩn mã nguồn) hoặc phân tích tổng hợp.
        \item Tác tử AI tự động thực hiện chu trình khám phá cây thư mục, tìm kiếm và đọc các tệp tin cần thiết.
    \end{itemize}
    \item \textbf{Phân hệ tra cứu và quản lý báo cáo}:
    \begin{itemize}
        \item Theo dõi trạng thái thực thi của tác vụ (\texttt{PENDING}, \texttt{PROCESSING}, \texttt{COMPLETED}, \texttt{ERROR}).
        \item Truy vấn kết quả đánh giá chi tiết dưới định dạng JSON có cấu trúc.
    \end{itemize}
\end{enumerate}

\subsection{Ma trận phân quyền người dùng (Role-Based Access Control)}
Bảng \ref{tab:rbac_matrix} mô tả phân quyền truy cập tài nguyên trong hệ thống theo vai trò người dùng.

\begin{table}[H]
\centering
\caption{Ma trận phân quyền truy cập tài nguyên (RBAC)}
\label{tab:rbac_matrix}
\renewcommand{\arraystretch}{1.2}
\begin{tabular}{|p{4.5cm}|p{3.5cm}|p{3.5cm}|}
\hline
\textbf{Tài nguyên / Thao tác} & \textbf{Khách (Chưa đăng nhập)} & \textbf{Người dùng (ROLE\_USER)} \\ \hline
Đăng ký / Đăng nhập & Cho phép & Cho phép \\ \hline
Xem/Tạo/Sửa/Xóa dự án & Từ chối (401 Unauthorized) & Cho phép (Chỉ dự án của chính mình) \\ \hline
Khởi tạo tác vụ phân tích AI & Từ chối (401 Unauthorized) & Cho phép (Chỉ trên dự án của chính mình) \\ \hline
Xem kết quả tác vụ & Từ chối (401 Unauthorized) & Cho phép (Kiểm tra quyền sở hữu IDOR) \\ \hline
\end{tabular}
\end{table}

\subsection{Yêu cầu phi chức năng}
\begin{itemize}
    \item \textbf{Tính bất đồng bộ}: Các tác vụ AI chạy nền để giải phóng luồng HTTP, giúp các điểm cuối API phản hồi kịp thời cho người dùng.
    \item \textbf{An toàn thông tin}: Kiểm tra quyền sở hữu tài nguyên để hạn chế truy cập chéo giữa các tài khoản và kiểm tra đường dẫn thư mục để ngăn ngừa duyệt tệp ngoài phạm vi dự án.
    \item \textbf{Tính linh hoạt của dữ liệu}: Sử dụng cơ sở dữ liệu quan hệ MySQL kết hợp trường dữ liệu JSON để lưu trữ linh hoạt cấu trúc báo cáo do tác tử AI sinh ra.
    \item \textbf{Khả năng đóng gói triển khai}: Hệ thống được cấu hình đóng gói qua Docker Compose nhằm thuận tiện cho việc thiết lập môi trường chạy thử nghiệm.
\end{itemize}

\section{Thiết kế kiến trúc tổng thể}
Hệ thống được thiết kế theo mô hình kiến trúc phân tầng (Layered Architecture), phân định rõ trách nhiệm giữa các thành phần. Sơ đồ kiến trúc tổng thể được thể hiện trong Hình \ref{fig:system_architecture}.

\begin{figure}[H]
\centering
\begin{tikzpicture}[
    node distance=1.1cm,
    layer/.style={rectangle, draw=black, fill=blue!5, thick, minimum width=13.5cm, minimum height=0.85cm, align=center, rounded corners=2pt, font=\small\bfseries},
    arrow/.style={-{Stealth[scale=1.0]}, thick}
]
    \node (client) [layer, fill=green!10] {Tầng Giao tiếp Khách hàng (Client: Postman / REST Client)};
    \node (security) [layer, fill=orange!10, below=0.7cm of client] {Tầng Điều khiển \& Bảo mật (Controller \& Security: JwtFilter, Auth, Project, AgentTask)};
    \node (service) [layer, fill=cyan!10, below=0.7cm of security] {Tầng Dịch vụ \& Điều phối (Service \& Async: ProjectService, AgentTaskService, TaskExecutor)};
    \node (agent) [layer, fill=purple!10, below=0.7cm of service] {Tầng Tác tử AI (Agent Core: CodeReviewAgent, Gemini API, FileSystemTools)};
    \node (data) [layer, fill=yellow!10, below=0.7cm of agent] {Tầng Dữ liệu (Persistence: Spring Data JPA Repositories, MySQL Database)};

    \draw [arrow] (client) -- node[right, font=\footnotesize] {HTTP REST (JSON / JWT)} (security);
    \draw [arrow] (security) -- node[right, font=\footnotesize] {Gọi Service} (service);
    \draw [arrow] (service) -- node[right, font=\footnotesize] {Kích hoạt @Async} (agent);
    \draw [arrow] (agent) -- node[right, font=\footnotesize] {Lưu JSON Kết quả} (data);
    \draw [arrow, dashed] (service.south east) -- ++(0.4,-0.35) |- (data.east);
\end{tikzpicture}
\caption{Sơ đồ kiến trúc phân tầng của hệ thống}
\label{fig:system_architecture}
\end{figure}

\section{Thiết kế cơ sở dữ liệu}
Hệ thống sử dụng cơ sở dữ liệu quan hệ MySQL gồm ba bảng thực thể: \texttt{users}, \texttt{projects} và \texttt{agent\_tasks}. Sơ đồ thực thể liên kết được thể hiện trong Hình \ref{fig:database_erd}.

\begin{figure}[H]
\centering
\begin{tikzpicture}[
    entity/.style={rectangle, draw=black, fill=blue!8, thick, align=left, font=\footnotesize, rounded corners=2pt, inner sep=4pt},
    rel/.style={-{Stealth[scale=1.0]}, thick}
]
    \node (user) [entity] at (0,0) {
        \begin{tabular}{l}
        \textbf{\large users} \\
        \hline
        \textbf{PK:} id (BIGINT) \\
        username (VARCHAR 50) \\
        email (VARCHAR 100) \\
        password (VARCHAR 255) \\
        role (VARCHAR 20) \\
        createdAt (DATETIME) \\
        updatedAt (DATETIME)
        \end{tabular}
    };

    \node (project) [entity] at (5.3,0) {
        \begin{tabular}{l}
        \textbf{\large projects} \\
        \hline
        \textbf{PK:} id (BIGINT) \\
        \textbf{FK:} user\_id (BIGINT) \\
        name (VARCHAR 100) \\
        description (VARCHAR 500) \\
        path (VARCHAR 500) \\
        language (VARCHAR 50) \\
        createdAt (DATETIME) \\
        updatedAt (DATETIME)
        \end{tabular}
    };

    \node (task) [entity] at (10.8,0) {
        \begin{tabular}{l}
        \textbf{\large agent\_tasks} \\
        \hline
        \textbf{PK:} id (BIGINT) \\
        \textbf{FK:} user\_id (BIGINT) \\
        \textbf{FK:} project\_id (BIGINT) \\
        type (VARCHAR 30) \\
        status (VARCHAR 20) \\
        input (JSON) \\
        output (JSON) \\
        errorMessage (VARCHAR 1000) \\
        createdAt (DATETIME) \\
        completedAt (DATETIME)
        \end{tabular}
    };

    \draw [rel] (user) -- node[above, font=\footnotesize] {1 : N} (project);
    \draw [rel] (project) -- node[above, font=\footnotesize] {1 : N} (task);
    \draw [rel] (user.south) -- ++(0,-0.75) -| node[pos=0.25, below, font=\footnotesize] {1 : N} (task.south);
\end{tikzpicture}
\caption{Sơ đồ thực thể liên kết (ERD) của hệ thống}
\label{fig:database_erd}
\end{figure}

\subsection{Từ điển dữ liệu chi tiết}
Bảng \ref{tab:data_dictionary} mô tả chi tiết các trường thuộc tính, kiểu dữ liệu và ràng buộc của các bảng trong cơ sở dữ liệu.

\begin{table}[H]
\centering
\caption{Đặc tả từ điển dữ liệu của hệ thống}
\label{tab:data_dictionary}
\renewcommand{\arraystretch}{1.2}
\begin{tabular}{|p{2.5cm}|p{2.5cm}|p{2.2cm}|p{2.2cm}|p{4.2cm}|}
\hline
\textbf{Tên bảng} & \textbf{Tên cột} & \textbf{Kiểu dữ liệu} & \textbf{Ràng buộc} & \textbf{Mô tả nghiệp vụ} \\ \hline
\texttt{users} & \texttt{id} & \texttt{BIGINT} & PK, Auto Inc & Khóa chính định danh tài khoản \\ \hline
\texttt{users} & \texttt{username} & \texttt{VARCHAR(50)} & Unique, Not Null & Tên đăng nhập người dùng \\ \hline
\texttt{users} & \texttt{email} & \texttt{VARCHAR(100)} & Unique, Not Null & Địa chỉ email \\ \hline
\texttt{users} & \texttt{password} & \texttt{VARCHAR(255)} & Not Null & Mật khẩu đã băm bằng BCrypt \\ \hline
\texttt{users} & \texttt{role} & \texttt{VARCHAR(20)} & Not Null & Vai trò (\texttt{ROLE\_USER}) \\ \hline
\texttt{projects} & \texttt{id} & \texttt{BIGINT} & PK, Auto Inc & Khóa chính định danh dự án \\ \hline
\texttt{projects} & \texttt{user\_id} & \texttt{BIGINT} & FK, Not Null & Khóa ngoại liên kết bảng \texttt{users} \\ \hline
\texttt{projects} & \texttt{name} & \texttt{VARCHAR(100)} & Not Null & Tên dự án mã nguồn \\ \hline
\texttt{projects} & \texttt{path} & \texttt{VARCHAR(500)} & Not Null & Đường dẫn thư mục mã nguồn \\ \hline
\texttt{agent\_tasks} & \texttt{id} & \texttt{BIGINT} & PK, Auto Inc & Khóa chính định danh tác vụ \\ \hline
\texttt{agent\_tasks} & \texttt{user\_id} & \texttt{BIGINT} & FK, Not Null & Khóa ngoại liên kết bảng \texttt{users} \\ \hline
\texttt{agent\_tasks} & \texttt{project\_id} & \texttt{BIGINT} & FK, Not Null & Khóa ngoại liên kết \texttt{projects} \\ \hline
\texttt{agent\_tasks} & \texttt{status} & \texttt{VARCHAR(20)} & Not Null & Trạng thái (\texttt{PENDING}, \dots) \\ \hline
\texttt{agent\_tasks} & \texttt{output} & \texttt{JSON} & Nullable & Báo cáo kết quả dạng JSON \\ \hline
\end{tabular}
\end{table}

\section{Thiết kế luồng xử lý đánh giá mã nguồn bất đồng bộ}

\subsection{Sơ đồ tuần tự thực thi tác vụ}
Luồng tương tác giữa các thành phần từ khi Client gửi yêu cầu đến khi lưu kết quả được mô tả trong Hình \ref{fig:review_sequence}.

\begin{figure}[H]
\centering
\begin{tikzpicture}[
    font=\footnotesize,
    lifeline/.style={draw=black!60, dashed, thick},
    actor/.style={rectangle, draw=black, fill=gray!15, thick, minimum width=1.7cm, minimum height=0.6cm, align=center, font=\footnotesize\bfseries, rounded corners=2pt},
    msg/.style={-{Stealth[scale=0.85]}, thick},
    reply/.style={-{Stealth[scale=0.85]}, dashed, thick}
]
    % Nodes
    \node (user) [actor] at (0,0) {Client};
    \node (ctrl) [actor] at (2.3,0) {Controller};
    \node (srv) [actor] at (4.7,0) {Service};
    \node (exec) [actor] at (7.3,0) {TaskExecutor};
    \node (agent) [actor] at (9.9,0) {AI Agent};
    \node (gemini) [actor] at (12.3,0) {Gemini};
    \node (db) [actor] at (14.5,0) {Database};

    % Lifelines
    \draw [lifeline] (user) -- (0,-6.6);
    \draw [lifeline] (ctrl) -- (2.3,-6.6);
    \draw [lifeline] (srv) -- (4.7,-6.6);
    \draw [lifeline] (exec) -- (7.3,-6.6);
    \draw [lifeline] (agent) -- (9.9,-6.6);
    \draw [lifeline] (gemini) -- (12.3,-6.6);
    \draw [lifeline] (db) -- (14.5,-6.6);

    % Messages
    \draw [msg] (0,-0.7) -- node[above, font=\scriptsize] {1. POST /api/tasks} (2.3,-0.7);
    \draw [msg] (2.3,-1.2) -- node[above, font=\scriptsize] {2. createTask()} (4.7,-1.2);
    \draw [msg] (4.7,-1.7) -- node[above, font=\scriptsize] {3. save(PENDING)} (14.5,-1.7);
    \draw [msg] (4.7,-2.2) -- node[above, font=\scriptsize] {4. executeAsync()} (7.3,-2.2);
    \draw [reply] (4.7,-2.7) -- node[above, font=\scriptsize] {5. TaskResponse} (2.3,-2.7);
    \draw [reply] (2.3,-3.2) -- node[above, font=\scriptsize] {6. 201 Created} (0,-3.2);
    \draw [msg] (7.3,-3.7) -- node[above, font=\scriptsize] {7. update(PROCESSING)} (14.5,-3.7);
    \draw [msg] (7.3,-4.2) -- node[above, font=\scriptsize] {8. reviewPath()} (9.9,-4.2);
    \draw [msg] (9.9,-4.7) -- node[above, font=\scriptsize] {9. Prompt + Tools} (12.3,-4.7);
    \draw [reply] (12.3,-5.1) -- node[above, font=\scriptsize] {10. Tool Call / Results} (9.9,-5.1);
    \draw [reply] (9.9,-5.6) -- node[above, font=\scriptsize] {11. JSON Output} (7.3,-5.6);
    \draw [msg] (7.3,-6.1) -- node[above, font=\scriptsize] {12. save(COMPLETED)} (14.5,-6.1);
\end{tikzpicture}
\caption{Sơ đồ tuần tự luồng xử lý đánh giá mã nguồn bất đồng bộ}
\label{fig:review_sequence}
\end{figure}

\subsection{Mô hình máy trạng thái của tác vụ (Task State Machine)}
Vòng đời của mỗi tác vụ phân tích mã nguồn được quản lý qua máy trạng thái hữu hạn, thể hiện trong Hình \ref{fig:task_state_machine}.

\begin{figure}[H]
\centering
\begin{tikzpicture}[
    node distance=2.6cm,
    state/.style={rectangle, draw=black, fill=blue!10, thick, minimum width=2.6cm, minimum height=0.9cm, align=center, rounded corners=6pt, font=\small\bfseries},
    arrow/.style={-{Stealth[scale=1.0]}, thick}
]
    \node (init) [circle, draw=black, fill=black, inner sep=2pt] {};
    \node (pending) [state, right=1.2cm of init] {PENDING};
    \node (processing) [state, fill=orange!15, right=of pending] {PROCESSING};
    \node (completed) [state, fill=green!15, right=of processing] {COMPLETED};
    \node (error) [state, fill=red!15, below=1.2cm of processing] {ERROR};

    \draw [arrow] (init) -- (pending);
    \draw [arrow] (pending) -- node[above, font=\footnotesize] {Nhận luồng xử lý} (processing);
    \draw [arrow] (processing) -- node[above, font=\footnotesize] {Hoàn thành} (completed);
    \draw [arrow] (processing) -- node[right, font=\footnotesize] {Lỗi / Ngoại lệ} (error);
    \draw [arrow] (pending) |- node[pos=0.75, above, font=\footnotesize] {Đường dẫn không hợp lệ} (error);
\end{tikzpicture}
\caption{Sơ đồ chuyển đổi trạng thái của tác vụ đánh giá}
\label{fig:task_state_machine}
\end{figure}

\section{Thiết kế giao diện lập trình ứng dụng (RESTful API)}
Hệ thống cung cấp các điểm cuối RESTful API phục vụ tương tác của người dùng, được tổng hợp trong Bảng \ref{tab:api_endpoints}.

\begin{table}[H]
\centering
\caption{Danh sách các điểm cuối API của hệ thống}
\label{tab:api_endpoints}
\renewcommand{\arraystretch}{1.2}
\begin{tabular}{|l|l|p{4.8cm}|l|l|}
\hline
\textbf{Method} & \textbf{Endpoint} & \textbf{Mô tả chức năng} & \textbf{Xác thực} & \textbf{Mã phản hồi} \\ \hline
POST & \texttt{/api/auth/register} & Đăng ký tài khoản người dùng mới & Công khai & 201 \\ \hline
POST & \texttt{/api/auth/login} & Đăng nhập và nhận mã thông báo JWT & Công khai & 200 \\ \hline
GET & \texttt{/api/projects} & Lấy danh sách dự án của người dùng & JWT & 200 \\ \hline
GET & \texttt{/api/projects/\{id\}} & Lấy thông tin chi tiết một dự án & JWT & 200 \\ \hline
POST & \texttt{/api/projects} & Khai báo dự án và đường dẫn mới & JWT & 201 \\ \hline
PUT & \texttt{/api/projects/\{id\}} & Cập nhật thông tin dự án & JWT & 200 \\ \hline
DELETE & \texttt{/api/projects/\{id\}} & Xóa dự án & JWT & 204 \\ \hline
POST & \texttt{/api/tasks} & Tạo tác vụ đánh giá mã nguồn & JWT & 201 \\ \hline
GET & \texttt{/api/tasks/\{id\}} & Tra cứu trạng thái và kết quả tác vụ & JWT & 200 \\ \hline
GET & \texttt{/api/tasks} & Lấy danh sách tất cả tác vụ của người dùng & JWT & 200 \\ \hline
GET & \texttt{/api/tasks/project/\{id\}} & Lấy danh sách tác vụ theo dự án & JWT & 200 \\ \hline
\end{tabular}
\end{table}

% ==============================================
% CHƯƠNG 3: QUÁ TRÌNH TRIỂN KHAI VÀ HIỆN THỰC HÓA HỆ THỐNG
% ==============================================
\chapter{QUÁ TRÌNH TRIỂN KHAI VÀ HIỆN THỰC HÓA HỆ THỐNG}

Công việc hiện thực hóa hệ thống được tổ chức qua 4 phân hệ: Xây dựng tác tử AI trên nền Spring AI; Thiết lập dịch vụ backend và xử lý bất đồng bộ; Cài đặt tầng bảo mật và kiểm soát truy cập hệ thống tệp; và Đóng gói triển khai với Docker.

\section{Tích hợp Spring AI và xây dựng tác tử đánh giá mã nguồn}

\subsection{Cấu hình ChatClient và tích hợp bộ công cụ tương tác hệ thống tệp}
Để mô hình ngôn ngữ lớn có thể đọc mã nguồn dự án trên máy chủ, tác tử được xây dựng dựa trên giao diện \texttt{ChatClient} của Spring AI và tích hợp bộ công cụ từ thư viện \texttt{spring-ai-agent-utils}:
\begin{itemize}
    \item \textbf{Cấu hình ChatClient}: Trong lớp \texttt{CodeReviewAgent}, tiến hành khởi tạo \texttt{ChatClient} với \texttt{ChatClient.Builder}, thiết lập lời nhắc hệ thống mặc định và đăng ký các công cụ thao tác tệp.
    \item \textbf{Đăng ký công cụ thao tác tệp}:
    \begin{itemize}
        \item \texttt{ListDirectoryTool}: Liệt kê tệp tin và thư mục con tại một đường dẫn, giúp tác tử nắm bắt cấu trúc phân gói của dự án.
        \item \texttt{GlobTool}: Tìm kiếm tệp tin dựa trên mẫu ký tự đại diện (như \texttt{**/*.java}, \texttt{**/pom.xml}) để định vị nhanh các tệp theo định dạng.
        \item \texttt{GrepTool}: Tìm kiếm từ khóa, tên phương thức hoặc chú thích mã nguồn trên phạm vi dự án mà không cần mở từng tệp.
        \item \texttt{FileSystemTools}: Cung cấp các thao tác đọc nội dung chi tiết của tệp mã nguồn phục vụ phân tích.
    \end{itemize}
\end{itemize}

\subsection{Thiết kế lời nhắc hệ thống (Prompt Engineering)}
Lời nhắc hệ thống đóng vai trò định hình hành vi và cấu trúc phản hồi của tác tử AI. Lời nhắc được lưu trong tệp \texttt{prompts/code-reviewer.md} và nạp qua đối tượng \texttt{Resource} của Spring:
\begin{itemize}
    \item \textbf{Nguyên tắc hoạt động an toàn}: Quy định tác tử chỉ hoạt động ở chế độ đọc (Read-only), không sửa đổi hay tạo mới tệp tin trong dự án.
    \item \textbf{Giới hạn phạm vi quét mã nguồn}: Hướng dẫn tác tử bỏ qua các thư mục biên dịch và cấu hình nội bộ (như \texttt{target}, \texttt{.git}, \texttt{.idea}, \texttt{node\_modules}) nhằm hạn chế lượng token gửi vào mô hình và giảm thời gian xử lý.
    \item \textbf{Quy trình phân tích tuần tự}: Hướng dẫn tác tử thực hiện theo các bước: (1) Khảo sát cây thư mục; (2) Đọc tệp cấu hình trung tâm (\texttt{pom.xml}, \texttt{application.yml}); (3) Đọc các tệp xử lý nghiệp vụ liên quan; (4) Tổng hợp báo cáo.
    \item \textbf{Ràng buộc định dạng phản hồi}: Yêu cầu mô hình trả về chuỗi JSON với cấu trúc gồm 3 phần: \texttt{summary} (tóm tắt tổng quan), \texttt{issues} (danh sách vấn đề phát hiện kèm mức độ, phân loại, đường dẫn tệp, mô tả và khuyến nghị khắc phục) và \texttt{suggestions} (khuyến nghị cải tiến).
\end{itemize}

\subsection{Cấu hình an toàn và tiền xử lý chuỗi phản hồi}
Trong quá trình thử nghiệm, mô hình Google Gemini có thể từ chối phân tích các đoạn mã chứa cấu hình bảo mật do cơ chế kiểm duyệt nội dung mặc định. Để xử lý vấn đề này:
\begin{itemize}
    \item Cấu hình tham số \texttt{GoogleGenAiSafetySetting} với ngưỡng \texttt{BLOCK\_ONLY\_HIGH} cho các danh mục kiểm duyệt trong cấu hình gọi API, giúp mô hình tiếp tục phân tích mã nguồn kỹ thuật mà không bị ngắt quãng.
    \item Xây dựng phương thức \texttt{cleanJsonResponse()} trong \texttt{CodeReviewAgent} nhằm loại bỏ các ký tự bọc định dạng Markdown (như \texttt{```json ... ```}) trước khi phân giải chuỗi kết quả sang đối tượng Java.
\end{itemize}

\section{Hiện thực hóa các dịch vụ backend và cơ chế xử lý bất đồng bộ}

\subsection{Kiến trúc phân tầng của hệ thống}
Hệ thống backend được tổ chức theo kiến trúc phân tầng:
\begin{itemize}
    \item \textbf{Tầng Controller}: Các lớp \texttt{AuthController}, \texttt{ProjectController} và \texttt{AgentTaskController} tiếp nhận yêu cầu HTTP, kiểm tra tính hợp lệ dữ liệu đầu vào (\texttt{@Valid}) và trả về phản hồi chuẩn RESTful.
    \item \textbf{Tầng Service}: \texttt{ProjectService} và \texttt{AgentTaskService} xử lý logic nghiệp vụ, xác thực quyền sở hữu tài nguyên và điều phối vòng đời tác vụ.
    \item \textbf{Tầng Repository}: Định nghĩa các interface \texttt{UserRepository}, \texttt{ProjectRepository} và \texttt{AgentTaskRepository} sử dụng Spring Data JPA để tương tác với cơ sở dữ liệu.
    \item \textbf{Tầng Entity và DTO}: Xây dựng các thực thể ánh xạ bảng cơ sở dữ liệu (\texttt{User}, \texttt{Project}, \texttt{AgentTask}) và các lớp DTO để truyền nhận dữ liệu với máy khách, tránh để lộ các thông tin nhạy cảm.
\end{itemize}

\subsection{Bộ quản lý luồng thực thi bất đồng bộ}
Quá trình phân tích của tác tử AI có thể kéo dài tùy theo quy mô của dự án. Để tránh việc phong tỏa luồng xử lý HTTP và gây quá thời gian chờ, cơ chế xử lý bất đồng bộ bằng \texttt{ThreadPoolTaskExecutor} được thiết lập như sau:
\begin{itemize}
    \item Trong lớp \texttt{AgentConfig}, cấu hình bean \texttt{agentThreadPoolExecutor} kiểu \texttt{ThreadPoolTaskExecutor} với các tham số:
    \begin{itemize}
        \item \texttt{corePoolSize = 5}: Duy trì 5 luồng xử lý thường trực.
        \item \texttt{maxPoolSize = 10}: Cho phép mở rộng tối đa 10 luồng khi có nhiều yêu cầu đồng thời.
        \item \texttt{queueCapacity = 50}: Hàng đợi chứa tối đa 50 tác vụ chờ.
        \item \texttt{threadNamePrefix = "AgentExecutor-"}: Tiền tố định danh luồng phục vụ theo dõi log.
    \end{itemize}
    \item Tách riêng logic thực thi tác tử AI sang lớp \texttt{AgentTaskExecutor} với phương thức \texttt{executeTaskAsync()} được đánh dấu \texttt{@Async("agentThreadPoolExecutor")}.
\end{itemize}

\subsection{Quản lý trạng thái và vòng đời tác vụ}
Vòng đời của tác vụ đánh giá mã nguồn được quản lý qua 4 trạng thái:
\begin{itemize}
    \item \textbf{Khởi tạo (\texttt{PENDING})}: Khi nhận yêu cầu phân tích, hệ thống lưu bản ghi tác vụ vào cơ sở dữ liệu với trạng thái \texttt{PENDING} và trả về mã tác vụ cho máy khách (mã 201 Created).
    \item \textbf{Đang xử lý (\texttt{PROCESSING})}: Luồng nền trong \texttt{AgentTaskExecutor} nhận tác vụ từ hàng đợi, chuyển trạng thái sang \texttt{PROCESSING} và bắt đầu gọi tác tử AI duyệt mã nguồn.
    \item \textbf{Hoàn thành (\texttt{COMPLETED})}: Khi tác tử AI hoàn thành phân tích và sinh báo cáo JSON hợp lệ, kết quả được lưu vào trường \texttt{output}, cập nhật trạng thái sang \texttt{COMPLETED} và ghi nhận thời điểm \texttt{completedAt}.
    \item \textbf{Lỗi (\texttt{ERROR})}: Nếu phát sinh lỗi trong quá trình thực thi (như lỗi kết nối API hoặc đường dẫn không hợp lệ), hệ thống ghi nhận thông điệp lỗi vào trường \texttt{errorMessage} và chuyển trạng thái sang \texttt{ERROR}.
\end{itemize}

\section{Bảo mật hệ thống và kiểm soát an toàn hệ thống tệp}

\subsection{Xác thực người dùng qua JWT và mã hóa mật khẩu}
\begin{itemize}
    \item Lớp \texttt{JwtTokenProvider} đảm nhiệm việc tạo lập, ký điện tử và kiểm tra mã thông báo JWT dựa trên thuật toán HMAC-SHA256.
    \item Lớp \texttt{JwtAuthenticationFilter} tích hợp vào chuỗi lọc của Spring Security, kiểm tra mã thông báo trong tiêu đề \texttt{Authorization: Bearer <token>} trước khi chuyển yêu cầu đến Controller.
    \item Mật khẩu người dùng được băm một chiều bằng \texttt{BCryptPasswordEncoder} trước khi lưu vào cơ sở dữ liệu.
\end{itemize}

\subsection{Kiểm soát quyền sở hữu tài nguyên (Hạn chế IDOR)}
Để hạn chế việc truy cập trái phép vào dự án hoặc tác vụ của tài khoản khác thông qua việc thay đổi định danh trên đường dẫn URL:
\begin{itemize}
    \item Trong \texttt{ProjectService} và \texttt{AgentTaskService}, hệ thống trích xuất định danh người dùng từ \texttt{SecurityContextHolder}.
    \item Đối chiếu định danh người dùng đang đăng nhập với trường \texttt{user\_id} của tài nguyên. Nếu không khớp, hệ thống từ chối yêu cầu và trả về lỗi 403 Forbidden hoặc 404 Not Found.
\end{itemize}

\subsection{Chuẩn hóa đường dẫn và kiểm tra an toàn thư mục (Hạn chế Path Traversal)}
Do tác tử AI cần truy cập hệ thống tệp trên máy chủ theo đường dẫn khai báo của dự án, hệ thống áp dụng cơ chế kiểm tra an toàn đường dẫn:
\begin{itemize}
    \item Trong lớp \texttt{AgentTaskExecutor}, phương thức \texttt{isPathAllowed()} thực hiện:
    \begin{enumerate}
        \item Chuẩn hóa đường dẫn đầu vào bằng \texttt{Path.normalize().toAbsolutePath()} để loại bỏ các đoạn duyệt ngược \texttt{..} và liên kết tượng trưng.
        \item So khớp đường dẫn đã chuẩn hóa với danh sách thư mục được cấp phép truy cập (\texttt{agent.allowed-paths}) khai báo trong \texttt{application.yml}.
    \end{enumerate}
    \item Nếu đường dẫn sau khi chuẩn hóa nằm ngoài phạm vi cho phép, hệ thống từ chối tác vụ, chuyển trạng thái sang \texttt{ERROR} và ghi nhận thông báo lỗi.
\end{itemize}

\section{Đóng gói hệ thống với Docker và cấu hình môi trường triển khai}

\subsection{Tệp cấu hình Dockerfile đa giai đoạn}
Nhằm tối ưu kích thước ảnh chứa ứng dụng:
\begin{itemize}
    \item \textbf{Giai đoạn biên dịch}: Sử dụng ảnh \texttt{eclipse-temurin:21-jdk-alpine} chứa JDK và Maven để biên dịch mã nguồn và đóng gói thành tệp \texttt{app.jar}.
    \item \textbf{Giai đoạn chạy ứng dụng}: Sử dụng ảnh gọn nhẹ \texttt{eclipse-temurin:21-jre-alpine} chứa môi trường JRE, sao chép tệp \texttt{app.jar} sang và cấu hình lệnh thực thi. Quá trình này giúp giảm dung lượng ảnh chạy từ hơn 600MB xuống còn khoảng 180MB.
\end{itemize}

\subsection{Cấu hình điều phối dịch vụ với Docker Compose}
Tệp \texttt{docker-compose.yml} được thiết lập để điều phối hai dịch vụ:
\begin{itemize}
    \item \textbf{Dịch vụ cơ sở dữ liệu (\texttt{mysql})}: Sử dụng MySQL 8.0, cấu hình ổ đĩa lưu trữ (\texttt{db\_data}) và kiểm tra trạng thái hoạt động (\texttt{healthcheck}) bằng lệnh \texttt{mysqladmin ping}.
    \item \textbf{Dịch vụ backend (\texttt{app})}: Phụ thuộc vào dịch vụ cơ sở dữ liệu với điều kiện \texttt{service\_healthy}, giúp ứng dụng backend khởi động sau khi MySQL đã sẵn sàng tiếp nhận kết nối.
    \item \textbf{Gắn kết thư mục mã nguồn (Volume)}: Gắn kết thư mục chứa các dự án mã nguồn vào đường dẫn \texttt{/workspace} trong container, cho phép tác tử AI đọc tệp tin theo phạm vi cấu hình.
\end{itemize}

% ==============================================
% CHƯƠNG 4: KẾT QUẢ THỬ NGHIỆM VÀ ĐÁNH GIÁ HỆ THỐNG
% ==============================================
\chapter{KẾT QUẢ THỬ NGHIỆM VÀ ĐÁNH GIÁ HỆ THỐNG}

\section{Kiểm thử chức năng và tích hợp hệ thống}

\subsection{Phương pháp và môi trường kiểm thử}
Để kiểm tra tính đúng đắn của các điểm cuối API, luồng xử lý dữ liệu và các cơ chế an toàn, hệ thống được kiểm thử chức năng thông qua công cụ Postman kết hợp theo dõi log máy chủ.

Các kịch bản kiểm thử được tổ chức theo 4 nhóm:
\begin{enumerate}
    \item \textbf{Nhóm xác thực và bảo mật}: Kiểm tra quá trình cấp phát, kiểm tra và từ chối mã thông báo JWT.
    \item \textbf{Nhóm phân quyền tài nguyên}: Kiểm tra việc cô lập dữ liệu giữa các tài khoản người dùng trên hệ thống (chống IDOR).
    \item \textbf{Nhóm kiểm tra an toàn đường dẫn và dữ liệu}: Kiểm tra khả năng từ chối các đường dẫn nằm ngoài phạm vi cho phép và xử lý chuỗi phản hồi từ mô hình AI.
    \item \textbf{Nhóm điều phối tác vụ bất đồng bộ}: Kiểm tra quá trình chuyển đổi trạng thái của tác vụ từ \texttt{PENDING} $\rightarrow$ \texttt{PROCESSING} $\rightarrow$ \texttt{COMPLETED}.
\end{enumerate}

\subsection{Chi tiết các ca kiểm thử}
Bảng \ref{tab:test_cases} tổng hợp 9 kịch bản kiểm thử đại diện đã thực hiện trên hệ thống backend.

\begin{table}[H]
\centering
\caption{Bảng kịch bản kiểm thử hệ thống trên Postman}
\label{tab:test_cases}
\renewcommand{\arraystretch}{1.2}
\begin{tabular}{|p{1.0cm}|p{3.2cm}|p{3.6cm}|p{3.8cm}|p{2.2cm}|}
\hline
\textbf{Mã} & \textbf{Chức năng} & \textbf{Dữ liệu đầu vào} & \textbf{Kết quả mong đợi} & \textbf{Đánh giá} \\ \hline
TC01 & Đăng ký tài khoản & Username và email hợp lệ & Tạo tài khoản, mã 201 Created kèm JWT & Đạt \\ \hline
TC02 & Đăng nhập sai & Tên đúng, mật khẩu sai & Trả về mã lỗi 401 Unauthorized & Đạt \\ \hline
TC03 & Thiếu mã xác thực & Request không kèm JWT header & Trả về mã lỗi 401 Unauthorized & Đạt \\ \hline
TC04 & Khởi tạo tác vụ & Dự án hợp lệ, thuộc user & Tạo task PENDING, trả về 201 Created & Đạt \\ \hline
TC05 & Kiểm tra IDOR & User A truy vấn task của User B & Bị từ chối, trả về mã 404/403 & Đạt \\ \hline
TC06 & Kiểm tra đường dẫn & Đường dẫn chứa \texttt{..} hoặc ngoài phạm vi & Bị từ chối, task chuyển sang ERROR & Đạt \\ \hline
TC07 & Định dạng JSON AI & LLM trả về JSON bọc markdown & Cắt bỏ markdown và parse thành công & Đạt \\ \hline
TC08 & Xử lý lỗi API & Khóa Gemini không hợp lệ / ngắt mạng & Bắt ngoại lệ, lưu errorMessage, status ERROR & Đạt \\ \hline
TC09 & Xử lý bất đồng bộ & Kích hoạt review dự án mẫu & Task chuyển PENDING $\rightarrow$ PROCESSING $\rightarrow$ COMPLETED & Đạt \\ \hline
\end{tabular}
\end{table}

\subsection{Nhận xét kết quả kiểm thử chức năng}
Trong phạm vi 9 kịch bản kiểm thử được thực hiện:
\begin{itemize}
    \item Hệ thống phản hồi các mã trạng thái HTTP tương ứng với từng tình huống nghiệp vụ.
    \item Cơ chế kiểm tra quyền sở hữu tài nguyên và kiểm tra an toàn đường dẫn đã từ chối các yêu cầu không hợp lệ.
    \item API tiếp nhận yêu cầu phân tích phản hồi thông tin tác vụ trong khi tiến trình AI được thực thi ở luồng nền.
\end{itemize}

\section{Kết quả đo độ phủ kiểm thử (JaCoCo)}

\subsection{Mục tiêu và phương pháp đo lường}
Để kiểm tra tính đúng đắn của logic nghiệp vụ và hỗ trợ phát hiện lỗi trong quá trình bảo trì, một bộ kiểm thử tự động đã được xây dựng. Bộ kiểm thử sử dụng khung JUnit 5, thư viện Mockito và công cụ đo độ phủ mã nguồn JaCoCo (\texttt{jacoco-maven-plugin} phiên bản 0.8.12).

Hai chỉ số đo lường chính bao gồm:
\begin{itemize}
    \item \textbf{Độ phủ dòng mã (Line Coverage)}: Tỷ lệ các dòng lệnh thực thi được kích hoạt bởi các ca kiểm thử.
    \item \textbf{Độ phủ nhánh điều kiện (Branch Coverage)}: Tỷ lệ các nhánh rẽ (\texttt{if-else}, toán tử ba ngôi) được duyệt qua các trường hợp \texttt{true} và \texttt{false}.
\end{itemize}

Cấu hình JaCoCo tập trung vào các lớp xử lý nghiệp vụ, điều khiển, tác tử và bảo mật; các lớp chỉ chứa cấu trúc dữ liệu thuần túy (Entity, DTO, Enum, Configuration) không được đưa vào diện tính toán độ phủ logic.

\subsection{Số liệu đo độ phủ thực tế}
Bộ kiểm thử bao gồm \textbf{50 bài kiểm thử tự động}, đạt kết quả thực thi \textbf{50/50 bài kiểm thử vượt qua (0 Failures, 0 Errors)}. Kết quả ghi nhận từ báo cáo JaCoCo trên toàn bộ dự án đạt:
\begin{itemize}
    \item \textbf{Độ phủ dòng mã}: \textbf{96.95\%} (350 / 361 dòng lệnh được kiểm thử, làm tròn 97.0\%).
    \item \textbf{Độ phủ nhánh điều kiện}: \textbf{89.86\%} (62 / 69 nhánh điều kiện được kiểm thử, làm tròn 89.9\%).
    \item \textbf{Độ phủ phương thức}: \textbf{100.0\%} (66 / 66 phương thức).
    \item \textbf{Độ phủ lớp nghiệp vụ}: \textbf{100.0\%} (11 / 11 lớp).
\end{itemize}

Bảng \ref{tab:test_coverage} trình bày chi tiết số liệu đo độ phủ theo từng lớp trong hệ thống. Đối với các lớp không chứa cấu trúc rẽ nhánh điều kiện (0/0 branch), chỉ số độ phủ nhánh được ghi nhận là \texttt{N/A}.

\begin{table}[H]
\centering
\caption{Kết quả đo độ phủ kiểm thử mã nguồn thực tế (JaCoCo)}
\label{tab:test_coverage}
\renewcommand{\arraystretch}{1.2}
\begin{tabular}{|p{3.2cm}|p{4.6cm}|p{3.0cm}|p{2.8cm}|}
\hline
\textbf{Phân hệ / Gói} & \textbf{Lớp (Class)} & \textbf{Độ phủ dòng} & \textbf{Độ phủ nhánh} \\ \hline
Controller & \texttt{AuthController} & 100.0\% (44/44) & 100.0\% (4/4) \\ \hline
Controller & \texttt{ProjectController} & 100.0\% (7/7) & N/A (0/0) \\ \hline
Controller & \texttt{AgentTaskController} & 100.0\% (5/5) & N/A (0/0) \\ \hline
Service & \texttt{ProjectService} & 100.0\% (50/50) & 100.0\% (2/2) \\ \hline
Service & \texttt{AgentTaskService} & 100.0\% (61/61) & 100.0\% (6/6) \\ \hline
Service & \texttt{AgentTaskExecutor} & 94.3\% (66/70) & 81.3\% (26/32) \\ \hline
Security & \texttt{JwtTokenProvider} & 93.8\% (30/32) & N/A (0/0) \\ \hline
Security & \texttt{JwtAuthenticationFilter} & 94.7\% (18/19) & 100.0\% (8/8) \\ \hline
Security & \texttt{CustomUserDetailsService} & 100.0\% (3/3) & N/A (0/0) \\ \hline
Security & \texttt{CustomUserDetails} & 100.0\% (7/7) & N/A (0/0) \\ \hline
Agent & \texttt{CodeReviewAgent} & 93.7\% (59/63) & 94.1\% (16/17) \\ \hline
\end{tabular}
\end{table}

\subsection{Phân tích kết quả kiểm thử theo từng phân hệ}
Dựa trên số liệu thu được từ báo cáo JaCoCo:

\begin{itemize}
    \item \textbf{Tầng Controller (56/56 dòng - 100\%)}:
    \begin{itemize}
        \item Toàn bộ các điểm cuối API trong \texttt{AuthController}, \texttt{ProjectController} và \texttt{AgentTaskController} đều được kiểm thử. 4 nhánh điều kiện trong \texttt{AuthController} đều được duyệt qua đầy đủ.
        \item Các bài kiểm thử bao phủ các tình huống: đăng ký tài khoản hợp lệ, từ chối khi trùng tên đăng nhập hoặc email; đăng nhập đúng cấp mã JWT, đăng nhập sai mật khẩu ném ngoại lệ \texttt{BadCredentialsException}; kiểm tra dữ liệu đầu vào và các mã phản hồi HTTP (\texttt{200 OK}, \texttt{201 Created}, \texttt{204 No Content}, \texttt{400 Bad Request}, \texttt{401 Unauthorized}).
    \end{itemize}

    \item \textbf{Tầng Service (177/181 dòng - 97.79\% Line, 34/40 nhánh - 85.00\% Branch)}:
    \begin{itemize}
        \item \texttt{ProjectService} và \texttt{AgentTaskService} đạt độ phủ 100\% ở cả dòng lệnh và nhánh rẽ. Các ca kiểm thử bao phủ việc kiểm tra quyền sở hữu tài nguyên, xử lý khi không tìm thấy dữ liệu (\texttt{ResourceNotFoundException}) và các thao tác thêm, sửa, xóa.
        \item \texttt{AgentTaskExecutor} đạt 94.29\% độ phủ dòng (66/70) và 81.25\% độ phủ nhánh (26/32). Các bài kiểm thử đã kích hoạt phương thức chuẩn hóa và kiểm tra đường dẫn \texttt{isPathAllowed()}, phân luồng tác vụ theo \texttt{TaskType} và bắt ngoại lệ khi xảy ra lỗi trong luồng ngầm. Một số dòng và nhánh chưa được bao phủ thuộc về các khối bắt ngoại lệ đường dẫn đặc biệt.
    \end{itemize}

    \item \textbf{Tầng Bảo mật (58/61 dòng - 95.08\% Line, 8/8 nhánh - 100\% Branch)}:
    \begin{itemize}
        \item 4 lớp bảo mật đạt độ phủ tổng thể 95.08\% dòng và 100\% nhánh rẽ.
        \item Các ca kiểm thử đã kiểm tra việc sinh mã JWT, trích xuất thông tin người dùng, kiểm tra mã hợp lệ, xử lý mã hết hạn (\texttt{ExpiredJwtException}), mã không đúng định dạng (\texttt{MalformedJwtException}) và lọc yêu cầu có hoặc không có tiêu đề \texttt{Authorization}.
    \end{itemize}

    \item \textbf{Tầng Tác tử AI (59/63 dòng - 93.65\% Line, 16/17 nhánh - 94.12\% Branch)}:
    \begin{itemize}
        \item Lớp \texttt{CodeReviewAgent} đạt 93.65\% độ phủ dòng và 94.12\% nhánh rẽ.
        \item Các bài kiểm thử đã kiểm tra việc khởi tạo lời nhắc theo các loại tác vụ (\texttt{SECURITY}, \texttt{PERFORMANCE}, \texttt{BUG}, \texttt{CODE\_QUALITY}) và kiểm tra phương thức \texttt{cleanJsonResponse()} trong việc loại bỏ khối định dạng Markdown.
    \end{itemize}
\end{itemize}

\subsection{Đánh giá về kết quả kiểm thử}
Kết quả kiểm thử cho thấy phần lớn các đoạn mã xử lý nghiệp vụ chính đã được kích hoạt thông qua các bài kiểm thử tự động, với độ phủ dòng đạt \textbf{96.95\%} và độ phủ nhánh đạt \textbf{89.86\%}. Mặc dù độ phủ kiểm thử cao không đồng nghĩa với việc loại bỏ hoàn toàn mọi lỗi tiềm ẩn trong thực tế, kết quả này cho thấy các luồng xử lý cơ bản, kiểm tra dữ liệu và xử lý ngoại lệ đã được kiểm tra độc lập.

\section{Số liệu thực nghiệm và đánh giá kết quả phân tích}

\subsection{Thiết lập môi trường thực nghiệm}
Quá trình thử nghiệm được thực hiện trên 5 dự án mã nguồn mẫu viết bằng Spring Boot (quy mô từ 15 đến 30 tệp Java, có cấu hình \texttt{pom.xml}, \texttt{application.yml} và các tầng Controller, Service, Repository). Hệ thống backend được chạy trong môi trường Docker với cấu hình 2 vCPU và 2GB RAM, kết nối với mô hình Google Gemini qua Spring AI.

\subsection{Số liệu thực nghiệm định lượng}
Bảng \ref{tab:empirical_metrics} thể hiện các chỉ số ghi nhận từ các lần chạy thử nghiệm của tác tử AI trên 5 dự án mẫu.

\begin{table}[H]
\centering
\caption{Số liệu thực nghiệm đánh giá mã nguồn của hệ thống}
\label{tab:empirical_metrics}
\renewcommand{\arraystretch}{1.2}
\begin{tabular}{|p{8.0cm}|p{6.0cm}|}
\hline
\textbf{Chỉ số thực nghiệm} & \textbf{Kết quả ghi nhận} \\ \hline
Quy mô dự án thử nghiệm trung bình & 22 tệp tin ($\approx$ 2.400 dòng mã) \\ \hline
Thời gian phân tích trung bình mỗi dự án & 28.5 giây (dao động 18s -- 42s) \\ \hline
Số lượt gọi công cụ (Tool Call) trung bình & 5.2 lượt / tác vụ \\ \hline
Tỷ lệ tệp tin cần đọc thực tế & 31.8\% (đọc trung bình 7/22 tệp trọng tâm) \\ \hline
Số lượng vấn đề phát hiện mức HIGH (Trung bình) & 1.4 vấn đề / dự án \\ \hline
Số lượng vấn đề phát hiện mức MEDIUM (Trung bình) & 3.2 vấn đề / dự án \\ \hline
Số lượng vấn đề phát hiện mức LOW (Trung bình) & 4.6 vấn đề / dự án \\ \hline
Mức tiêu thụ bộ nhớ RAM máy chủ backend & 420 MB -- 510 MB \\ \hline
Tỷ lệ phản hồi cấu trúc JSON hợp lệ & 100\% (sau khi làm sạch Markdown) \\ \hline
\end{tabular}
\end{table}

\subsection{Phân tích kết quả thực nghiệm}
Từ các số liệu ghi nhận trong Bảng \ref{tab:empirical_metrics}:
\begin{itemize}
    \item \textbf{Cơ chế gọi công cụ}: Thông qua \texttt{ListDirectoryTool} và \texttt{GrepTool}, tác tử AI đọc nội dung của khoảng 31.8\% tổng số tệp tin trong dự án (chủ yếu là tệp cấu hình, thực thể và lớp dịch vụ chính). Cách tiếp cận này giúp hạn chế gửi các tệp tin không cần thiết vào mô hình ngôn ngữ lớn so với phương pháp nạp toàn bộ mã nguồn.
    \item \textbf{Khả năng phát hiện vấn đề}: Trên các dự án mẫu thử nghiệm, tác tử AI đã ghi nhận một số vấn đề:
    \begin{itemize}
        \item \textit{Bảo mật}: Phát hiện khóa bí mật để mặc định trong tệp cấu hình, thiếu kiểm tra quyền truy cập tài nguyên.
        \item \textit{Hiệu năng}: Cảnh báo truy vấn cơ sở dữ liệu thiếu chỉ mục hoặc vòng lặp chưa tối ưu.
        \item \textit{Quy chuẩn mã nguồn}: Đề xuất bổ sung kiểm thử đơn vị, tách nhỏ phương thức dài và chuẩn hóa xử lý ngoại lệ.
    \end{itemize}
    \item \textbf{Sử dụng tài nguyên}: Trong phạm vi các lần thử nghiệm, mức tiêu thụ bộ nhớ RAM của ứng dụng backend dao động từ 420 MB đến 510 MB và không ghi nhận xu hướng tăng bất thường.
\end{itemize}

\subsection{Minh họa kết quả đánh giá mã nguồn thực tế}
Kết quả đánh giá mã nguồn thực tế do tác tử AI sinh ra sau khi hoàn thành tác vụ được minh họa thông qua phản hồi JSON của điểm cuối API \texttt{GET /api/tasks/1} trên giao diện công cụ Postman như trong Hình \ref{fig:sample_json_result}.

\begin{figure}[H]
\centering
\includegraphics[width=0.95\textwidth]{postman_task_result.png}
\caption{Minh họa kết quả truy vấn tác vụ đánh giá mã nguồn trả về định dạng JSON trên Postman}
\label{fig:sample_json_result}
\end{figure}

Thông qua kết quả trả về từ API trong Hình \ref{fig:sample_json_result}, có thể rút ra một số nhận định kỹ thuật về quá trình thực thi của hệ thống:
\begin{itemize}
    \item \textbf{Vòng đời xử lý bất đồng bộ hoàn chỉnh}: Tác vụ được hệ thống tiếp nhận và bắt đầu thực thi ngầm lúc \texttt{22:04:33}, hoàn thành lúc \texttt{22:04:49} với tổng thời gian xử lý khoảng 15.3 giây. Trạng thái tác vụ chuyển đổi thành công sang \texttt{COMPLETED} và dữ liệu kết quả được lưu trữ đầy đủ trong cơ sở dữ liệu MySQL mà không làm nghẽn luồng tiếp nhận yêu cầu HTTP.
    \item \textbf{Cấu trúc dữ liệu đầu ra chuẩn hóa}: Kết quả phản hồi tuân thủ cấu trúc JSON phân cấp gồm các phần thông tin tổng quát (\texttt{summary}), danh sách các vấn đề phát hiện (\texttt{issues}) và đề xuất cải tiến (\texttt{suggestions}), thuận tiện cho việc trích xuất và hiển thị trên giao diện người dùng.
    \item \textbf{Khả năng hiểu ngữ cảnh và khuyến nghị thực tế}: Tác tử AI đã tự động phân tích cấu trúc mã nguồn phân tầng của dự án (Spring Boot, Spring Security, JWT stateless, Thread Pool), đồng thời đưa ra các đề xuất tối ưu kỹ thuật phù hợp với thực tế phát triển phần mềm như bổ sung bộ nhớ đệm (Cache), tinh chỉnh connection pool (HikariCP) và tích hợp tài liệu giao diện OpenAPI/Swagger.
\end{itemize}

\section{Các vấn đề kỹ thuật phát sinh và biện pháp khắc phục}
Trong quá trình phát triển và thử nghiệm hệ thống, một số vấn đề kỹ thuật đã phát sinh và được xử lý:
\begin{itemize}
    \item \textbf{Quá thời gian chờ kết nối HTTP}:
    \begin{itemize}
        \item \textit{Nguyên nhân}: Ban đầu hệ thống thực thi tác tử AI đồng bộ trong Controller. Khi phân tích các dự án lớn, thời gian xử lý kéo dài gây gián đoạn kết nối HTTP.
        \item \textit{Cách khắc phục}: Chuyển sang mô hình bất đồng bộ với \texttt{ThreadPoolTaskExecutor} và phương thức \texttt{@Async}, lưu trạng thái tác vụ vào cơ sở dữ liệu để người dùng truy vấn sau.
    \end{itemize}
    \item \textbf{Lỗi phân giải JSON do định dạng Markdown}:
    \begin{itemize}
        \item \textit{Nguyên nhân}: Mô hình đôi khi tự động bọc chuỗi kết quả trong khối \texttt{```json ... ```}, dẫn đến lỗi khi chuyển đổi sang đối tượng Java.
        \item \textit{Cách khắc phục}: Bổ sung phương thức \texttt{cleanJsonResponse()} trong \texttt{CodeReviewAgent} để cắt bỏ các ký tự định dạng trước khi phân giải chuỗi.
    \end{itemize}
    \item \textbf{Kiểm soát truy xuất tệp}:
    \begin{itemize}
        \item \textit{Nguyên nhân}: Đường dẫn dự án có thể chứa ký tự điều hướng ngược (\texttt{../}) để truy cập ngoài phạm vi thư mục được phân công.
        \item \textit{Cách khắc phục}: Cài đặt phương thức \texttt{isPathAllowed()} trong \texttt{AgentTaskExecutor} để chuẩn hóa đường dẫn tuyệt đối và đối chiếu với danh sách thư mục được cấp phép trước khi thực thi.
    \end{itemize}
\end{itemize}

\section{Đánh giá tổng kết hệ thống}

\subsection{Kết quả đạt được}
\begin{itemize}
    \item Tác tử AI có khả năng tự động khám phá cấu trúc thư mục và đọc tệp mã nguồn thông qua cơ chế gọi công cụ.
    \item Phân tách luồng xử lý AI khỏi luồng HTTP giúp hệ thống duy trì khả năng tiếp nhận yêu cầu từ máy khách.
    \item Tích hợp xác thực JWT, kiểm tra quyền sở hữu tài nguyên (hạn chế IDOR) và kiểm tra an toàn đường dẫn (hạn chế Path Traversal).
    \item Kết quả đánh giá được lưu trữ dưới định dạng JSON có cấu trúc, thuận tiện cho việc truy vấn và tích hợp.
    \item Ứng dụng và cơ sở dữ liệu được đóng gói qua Docker Compose, thuận lợi cho việc tái lập môi trường chạy thử nghiệm.
\end{itemize}

\subsection{Hạn chế tồn tại}
\begin{itemize}
    \item Hệ thống hiện tại chỉ cung cấp dịch vụ backend và giao diện RESTful API, chưa có giao diện web trực quan cho người dùng cuối.
    \item Chưa có cơ chế thông báo theo thời gian thực (như WebSocket hay Server-Sent Events); người dùng cần truy vấn định kỳ (Polling) để xem trạng thái tác vụ.
    \item Tốc độ và độ ổn định của quá trình phân tích phụ thuộc vào kết nối mạng và giới hạn tần suất yêu cầu (Rate Limit) của dịch vụ API Google Gemini bên ngoài.
\end{itemize}

% ==============================================
% CHƯƠNG 5: KẾT LUẬN VÀ HƯỚNG PHÁT TRIỂN
% ==============================================
\chapter{KẾT LUẬN VÀ HƯỚNG PHÁT TRIỂN}

\section{Tổng kết các kết quả đạt được}
Trong khuôn khổ đợt thực tập tốt nghiệp, đề tài ``Nghiên cứu và xây dựng AI Agent hỗ trợ Code Review sử dụng Spring AI'' đã hoàn thành các mục tiêu kỹ thuật đề ra:

\subsection{Về mặt nghiên cứu và tiếp cận công nghệ}
\begin{itemize}
    \item \textbf{Tìm hiểu nguyên lý tác tử AI và cơ chế gọi công cụ}: Nắm được nguyên lý hoạt động của mô hình suy luận nhiều bước ReAct và cơ chế gọi công cụ ngoại vi (Tool Calling) để tương tác có kiểm soát với hệ thống tệp.
    \item \textbf{Ứng dụng khung làm việc Spring AI}: Sử dụng Spring AI kết nối với mô hình Google Gemini, áp dụng giao diện \texttt{ChatClient}, nạp lời nhắc từ tệp tài nguyên và đăng ký các công cụ thao tác tệp.
    \item \textbf{Thiết kế lời nhắc hệ thống}: Xây dựng tệp lời nhắc \texttt{prompts/code-reviewer.md} hướng dẫn mô hình thực hiện rà soát theo quy trình, loại bỏ các thư mục không cần thiết và định dạng kết quả trả về dưới dạng JSON có cấu trúc.
\end{itemize}

\subsection{Về mặt hiện thực hóa hệ thống}
\begin{itemize}
    \item \textbf{Xây dựng kiến trúc phân tầng}: Cài đặt hoàn chỉnh hệ thống backend trên nền tảng Java 21 và Spring Boot theo mô hình phân tầng (Controller, Service, Repository, DTO).
    \item \textbf{Xử lý tác vụ bất đồng bộ}: Quản lý vòng đời tác vụ qua 4 trạng thái (\texttt{PENDING}, \texttt{PROCESSING}, \texttt{COMPLETED}, \texttt{ERROR}) kết hợp với \texttt{ThreadPoolTaskExecutor}, giúp giải phóng luồng HTTP và trả về phản hồi kịp thời cho người dùng.
    \item \textbf{Cài đặt các cơ chế bảo mật cơ bản}: Áp dụng xác thực không trạng thái JWT, băm mật khẩu bằng BCrypt, kiểm tra quyền sở hữu tài nguyên nhằm hạn chế IDOR và kiểm tra an toàn đường dẫn với \texttt{isPathAllowed()} nhằm hạn chế Path Traversal.
    \item \textbf{Lưu trữ dữ liệu kết hợp}: Sử dụng cơ sở dữ liệu quan hệ MySQL kết hợp trường dữ liệu JSON (\texttt{@JdbcTypeCode(SqlTypes.JSON)}) để lưu trữ linh hoạt báo cáo do tác tử AI sinh ra.
\end{itemize}

\subsection{Về mặt kiểm thử và đóng gói triển khai}
\begin{itemize}
    \item \textbf{Kiểm thử tự động}: Xây dựng 50 bài kiểm thử đơn vị và tích hợp với JUnit 5 và Mockito, đạt độ phủ dòng \textbf{96.95\%} và độ phủ nhánh \textbf{89.86\%} theo kết quả đo lường từ JaCoCo.
    \item \textbf{Đóng gói với Docker}: Xây dựng tệp \texttt{Dockerfile} đa giai đoạn giúp giảm kích thước ảnh chạy xuống khoảng 180MB và tệp \texttt{docker-compose.yml} điều phối dịch vụ backend cùng cơ sở dữ liệu MySQL.
\end{itemize}

\section{Đánh giá giá trị thực tiễn và tính ứng dụng}
Hệ thống được xây dựng mang lại một số giá trị trong việc hỗ trợ quy trình phát triển phần mềm:
\begin{itemize}
    \item \textbf{Hỗ trợ sơ lọc mã nguồn}: Tác tử AI có thể rà soát các vấn đề phổ biến như khóa bí mật để mặc định trong tệp cấu hình hoặc các đoạn truy vấn chưa tối ưu, hỗ trợ lập trình viên kiểm tra trước khi đưa vào quy trình đánh giá chính thức.
    \item \textbf{Bổ trợ cho các công cụ phân tích tĩnh}: Kết hợp việc phân tích mã nguồn dựa trên mô hình ngôn ngữ lớn để đưa ra lời giải thích và khuyến nghị bằng ngôn ngữ tự nhiên bên cạnh các quy tắc kiểm tra cố định.
    \item \textbf{Tối ưu phạm vi đọc tệp}: Thông qua cơ chế gọi công cụ, tác tử chỉ đọc nội dung các tệp tin liên quan (chiếm khoảng 31.8\% số tệp trong các lần thử nghiệm), hạn chế gửi toàn bộ dự án vào mô hình.
\end{itemize}

\section{Bài học kinh nghiệm trong quá trình thực tập}
Quá trình thực hiện đề tài giúp sinh viên tích lũy thêm các kiến thức và kỹ năng thực tế:
\begin{itemize}
    \item \textbf{Kỹ năng tích hợp AI vào ứng dụng}: Nắm vững cách xây dựng công cụ ngoại vi và thiết kế cấu trúc dữ liệu trao đổi giữa mô hình ngôn ngữ lớn và hệ thống backend.
    \item \textbf{Kỹ năng thiết kế xử lý bất đồng bộ}: Hiểu rõ cách quản lý luồng nền, xử lý quá thời gian chờ và đồng bộ hóa trạng thái dữ liệu trong Spring Boot.
    \item \textbf{Nhận thức về an toàn thông tin}: Hiểu được tầm quan trọng của việc kiểm tra đường dẫn và kiểm soát quyền truy cập tài nguyên khi cho phép hệ thống tương tác với tệp tin trên máy chủ.
    \item \textbf{Kỹ năng kiểm thử mã nguồn}: Làm quen với việc sử dụng Mockito để cô lập phụ thuộc và đọc báo cáo JaCoCo để cải thiện chất lượng kiểm thử.
\end{itemize}

\section{Hướng phát triển trong tương lai}
Để hoàn thiện hệ thống, các hướng phát triển tiếp theo bao gồm:

\subsection{Xây dựng giao diện người dùng (Frontend Dashboard)}
Phát triển giao diện web (sử dụng React hoặc Next.js) phục vụ quản lý dự án, theo dõi tiến độ phân tích và hiển thị kết quả đánh giá trực quan kèm tính năng tô sáng cú pháp mã nguồn.

\subsection{Tích hợp cơ chế thông báo theo thời gian thực}
Sử dụng giao thức WebSocket hoặc Server-Sent Events (SSE) để máy chủ tự động đẩy thông báo tiến độ và kết quả phân tích đến giao diện người dùng ngay khi tác vụ hoàn thành, thay cho cơ chế truy vấn định kỳ hiện tại.

\subsection{Mở rộng hệ thống đa tác tử (Multi-Agent System)}
Nghiên cứu mở rộng thêm các tác tử chuyên trách như tác tử đề xuất bản vá mã nguồn, tác tử giải thích luồng hoạt động của module và tác tử hỗ trợ sinh mã kiểm thử đơn vị.

\subsection{Tích hợp vào chu trình CI/CD}
Xây dựng tiện ích tích hợp vào GitHub Actions hoặc GitLab CI để tự động kích hoạt tác vụ phân tích mỗi khi có yêu cầu gộp mã (Pull Request) mới.

% ==============================================
% TÀI LIỆU THAM KHẢO
% ==============================================
\phantomsection
\chapter*{Tài liệu tham khảo}
\addcontentsline{toc}{chapter}{Tài liệu tham khảo}

\begin{enumerate}
    \item Spring Boot Reference Documentation: \url{https://spring.io/projects/spring-boot}
    \item Spring AI Reference Documentation: \url{https://docs.spring.io/spring-ai/reference/}
    \item Google Gemini API Documentation: \url{https://ai.google.dev/docs}
    \item OpenAI Function Calling Guide: \url{https://platform.openai.com/docs/guides/function-calling}
    \item Anthropic Tool Use Overview: \url{https://docs.anthropic.com/en/docs/build-with-claude/tool-use}
    \item OWASP API Security Top 10 (Broken Object Level Authorization): \url{https://owasp.org/API-Security/editions/2023/en/0xa1-broken-object-level-authorization/}
    \item OWASP Path Traversal Prevention Cheat Sheet: \url{https://owasp.org/www-community/attacks/Path_Traversal}
    \item JaCoCo Java Code Coverage Library Documentation: \url{https://www.eclemma.org/jacoco/trunk/doc/index.html}
    \item Docker Official Documentation (Get Started): \url{https://docs.docker.com/get-started/}
    \item GitHub Documentation (About Git): \url{https://docs.github.com/en/get-started/using-git/about-git}
    \item Devteria Identity Service (Spring Boot \& JWT Authentication): \url{https://github.com/devteria/identity-service}
    \item Devteria Spring AI Demo (LLM Integration with Spring AI): \url{https://github.com/devteria/spring-ai-demo}
    \item Devteria Data Encryption Demo (Data Security in Spring Boot): \url{https://github.com/devteria/data-encryption-demo}
\end{enumerate}

\end{document}
